%% Expanded sections: Genetic Architecture & Therapeutic Landscape
%% Raw LaTeX content — no preamble, no \begin{document}


%% ====================================================================
%%  SECTION: GENETIC ARCHITECTURE
%% ====================================================================

\sectitle{Genetic Architecture of Emotional Dysregulation}

Clinicians have long observed that emotional dysregulation clusters in families,
but only in the past decade have researchers mapped its polygenic architecture
with sufficient statistical power to distinguish signal from noise.\cit{1}
Three converging lines of genome-wide evidence now delineate the heritable
basis of dysregulated affect: large-scale GWAS of neuroticism, depression, and
anxiety; functional enrichment analyses that situate risk loci within
neurotransmitter and neuroendocrine pathways; and epigenetic studies that reveal
how environmental adversity inscribes itself onto these same genes.

\subsection*{Genome-Wide Association: Scale and Convergence}

Neuroticism — the dispositional tendency toward negative affect that most
personality theorists regard as the temperamental substrate of emotional
dysregulation — has attracted at least 20 GWAS to date.\cit{2}
The largest of these, GCST006476, genotyped 449,484 participants and identified
136 independent genomic loci, while a subsequent study (GCST007339) extended
discovery to 523,783 individuals with an additional 59,206-person replication
cohort.\cit{3,4}
These studies transformed neuroticism from a ``soft'' personality construct into
a measurable polygenic trait whose genetic signal rivals that of height or
BMI in magnitude, if not in per-locus effect size.

Depression has generated an even larger literature — 27 GWAS and counting.
The most recent and powerful, GCST90319327, assembled 371,184 cases against
978,703 controls, making it the largest depression GWAS ever
conducted.\cit{5}
This unprecedented sample yielded hundreds of genome-wide significant loci and,
crucially, demonstrated that many of these loci overlap with neuroticism risk
variants — a pattern consistent with the hypothesis that neuroticism and
depression share a common heritable core rather than merely co-occurring
phenotypically.

Anxiety, though less extensively studied at the genomic level, has nonetheless
accumulated 11 GWAS.\cit{6}
The largest, GCST90468110, drew on 394,626 participants and identified multiple
loci that again overlapped with both neuroticism and depression signals.
A critical methodological feature strengthens these convergence claims: the UK
Biobank cohort ($N \approx 394{,}626$) has contributed GWAS data for all three
phenotypes — neuroticism, anxiety, and depression — enabling researchers to
estimate genetic correlations within the same participants rather than across
cohorts with different ascertainment biases.\cit{7}
The resulting genetic correlation between neuroticism and depression
($r_g \approx 0.80$) confirms that these phenotypes share the vast majority of
their common-variant architecture.\cit{8}
This overlap does not mean the phenotypes are identical; rather, it suggests
that a shared genetic liability disposes individuals toward emotional
dysregulation, which then manifests as neuroticism, depression, or anxiety
depending on developmental context and environmental exposure.

\subsection*{From Risk Loci to Functional Pathways}

Identifying hundreds of risk loci accomplishes little unless researchers can
trace those loci to biological mechanisms.
To bridge this gap, we performed protein--protein interaction and functional
enrichment analyses using the STRING database on eight candidate genes
repeatedly implicated across emotional dysregulation phenotypes:
\gene{SLC6A4}, \gene{MAOA}, \gene{COMT}, \gene{BDNF}, \gene{NR3C1},
\gene{CRH}, \gene{FKBP5}, and \gene{OXTR}.\cit{9}

Gene Ontology (GO) Biological Process enrichment revealed that these eight
genes converge on a surprisingly coherent set of functions.
Five of the eight annotate to ``behavior'' (FDR\,=\,0.011), and four to
``learning or memory'' (FDR\,=\,0.013) — confirming that the genetic
architecture of emotional dysregulation is not diffusely distributed across the
genome but concentrated in pathways that govern affective and cognitive
processing.\cit{10}
More specific enrichments illuminate the molecular mechanisms through which
these genes operate.
\gene{MAOA} and \gene{COMT} jointly enrich for ``neurotransmitter catabolic
process'' (FDR\,=\,0.021), reflecting their shared role in degrading
monoamine neurotransmitters — the same molecules targeted by first-line
antidepressants.
\gene{NR3C1} and \gene{CRH} enrich for ``glucocorticoid metabolic process''
(FDR\,=\,0.021), placing the hypothalamic--pituitary--adrenal (HPA) axis at
the center of genetic risk.
\gene{NR3C1} and \gene{OXTR} enrich for ``maternal behavior''
(FDR\,=\,0.021), a finding that resonates with developmental evidence linking
early caregiving quality to long-term emotion regulation capacity.\cit{11}
Finally, \gene{SLC6A4}, \gene{MAOA}, and \gene{COMT} together enrich for
``regulation of neurotransmitter levels'' (FDR\,=\,0.048), underscoring that
synaptic monoamine homeostasis represents a convergent node of genetic risk.

KEGG pathway enrichment extended these findings to disease-relevant circuitry.
The strongest enrichment appeared for ``Alcoholism''
(\gene{CRH}+\gene{MAOA}+\gene{BDNF}; FDR\,=\,0.008), consistent with
extensive comorbidity between emotional dysregulation and substance use
disorders.\cit{12}
``Tyrosine metabolism'' (\gene{MAOA}+\gene{COMT}; FDR\,=\,0.016) and
``Cocaine addiction'' (\gene{MAOA}+\gene{BDNF}; FDR\,=\,0.020) further
implicate catecholamine pathways, while ``Neuroactive ligand--receptor
interaction'' (\gene{NR3C1}+\gene{CRH}+\gene{OXTR}; FDR\,=\,0.021) captures
the neuropeptide and steroid signaling axis.
Taken together, these enrichments delineate two principal molecular systems
through which genetic variation shapes emotional dysregulation: monoamine
neurotransmission and neuroendocrine stress signaling.

\subsection*{Molecular Mechanisms at Single-Gene Resolution}

The enrichment analyses identify pathways; understanding how individual genes
within those pathways contribute to dysregulation requires examining their
molecular functions in detail.

\gene{SLC6A4} (UniProt P31645) encodes the sodium-dependent serotonin
transporter, which clears serotonin from the synaptic cleft via Na\textsuperscript{+}/K\textsuperscript{+}
cotransport.\cit{13}
This transporter is essential for serotonin homeostasis throughout the central
nervous system; its dysfunction disrupts the Reactome-annotated ``serotonin
clearance from the synaptic cleft'' pathway (R-HSA-380615) and, more broadly,
``SLC-mediated neurotransmitter transport'' (R-HSA-442660).\cit{14}
Allen Brain Atlas data confirm that \gene{Slc6a4} expression in the mouse
(Entrez 15567, homologene 817) concentrates in the raphe nuclei — the brain's
primary serotonin source — but also extends to cortical and limbic
regions that subserve emotion regulation.\cit{15}
Functional polymorphisms in \gene{SLC6A4}, particularly the 5-HTTLPR
insertion/deletion variant, alter transporter expression and have been
associated with amygdala reactivity to threat, although effect sizes in
candidate-gene studies remain contentious in the post-GWAS era.\cit{16}

\gene{FKBP5} (UniProt Q13451) operates through an entirely different molecular
logic.
Rather than modulating neurotransmitter levels directly, FKBP5 participates in
the HSP90 chaperone cycle for steroid hormone receptors
(Reactome R-HSA-3371497), where it acts as a co-chaperone that reduces
glucocorticoid receptor (GR) sensitivity.\cit{17}
When cortisol binds GR, FKBP5 is displaced by FKBP4, allowing the
receptor--ligand complex to translocate to the nucleus and activate
transcription.
Critically, one of the genes that GR activation upregulates is \gene{FKBP5}
itself, creating an ultrashort negative feedback loop: more cortisol produces
more FKBP5, which dampens GR sensitivity, which limits further cortisol-driven
transcription.\cit{18}
Polymorphisms in \gene{FKBP5} that increase its expression therefore blunt HPA
axis negative feedback, prolonging cortisol exposure after stress — a mechanism
that converges with the \gene{NR3C1}--\gene{CRH} glucocorticoid pathway
identified in our enrichment analysis.

These two genes exemplify the broader principle that genetic risk for emotional
dysregulation distributes across two interacting systems: fast synaptic
signaling (serotonin, catecholamines) and slow neuroendocrine modulation
(cortisol, oxytocin).
KEGG pathway annotations reinforce this duality — serotonin receptor signaling
(map07211) and oxytocin signaling (map04921) represent parallel but interacting
cascades, while cortisol synthesis and secretion (map04927) provides the
hormonal backdrop against which both operate.\cit{19}

\subsection*{Epigenetic Mediation: Where Genes Meet Environment}

If the GWAS literature identifies \emph{which} genes confer risk, epigenetic
research reveals \emph{how} environmental adversity activates that risk.
Udalov and colleagues (2026) demonstrated that early-life stress produces
measurable DNA methylation changes at three of the genes identified in our
enrichment analysis — \gene{NR3C1}, \gene{FKBP5}, and \gene{BDNF} — and that
these methylation signatures are partially reversible with
intervention.\cit{20}
This partial reversibility carries profound clinical implications: it means that
the epigenetic marks left by adversity are not permanent sentences but
modifiable targets.

MicroRNA profiling has added a post-transcriptional layer to this picture.
miR-16, miR-124, miR-132, miR-135a, and miR-34c — all associated with limbic
circuit remodeling — regulate the translation of mRNAs from genes in both the
monoamine and neuroendocrine pathways.\cit{21}
miR-16, for example, targets \gene{SLC6A4} mRNA and thereby modulates
serotonin transporter density, while miR-135a regulates both \gene{SLC6A4} and
the serotonin receptor \gene{HTR1A}.\cit{22}
These miRNAs respond to both stress exposure and therapeutic intervention,
suggesting that they function as molecular mediators through which experience
reshapes the neural circuits of emotion regulation.

The emerging picture, then, is one of layered vulnerability: common genetic
variants establish a baseline sensitivity to emotional dysregulation; epigenetic
modifications — driven by early adversity, chronic stress, or their absence —
tune gene expression within those variants; and post-transcriptional regulation
by miRNAs provides a rapid, reversible adjustment mechanism that tracks ongoing
environmental conditions.
This architecture is not merely descriptive.
It generates testable predictions about which individuals will respond to which
treatments — predictions that the next section begins to evaluate.


%% ====================================================================
%%  TRANSITION: FROM ARCHITECTURE TO INTERVENTION
%% ====================================================================

\bigskip

The molecular architecture described above does more than catalog risk factors;
it specifies the biological systems that treatments must engage to restore
regulatory capacity.
When genetic and epigenetic evidence converges on monoamine neurotransmission,
clinicians gain a mechanistic rationale for serotonergic pharmacotherapy — not
as a blunt correction of ``chemical imbalance,'' but as a targeted modulation
of the \gene{SLC6A4}--dependent clearance pathway that GWAS, Reactome, and
Allen Brain Atlas data independently implicate.\cit{23}
When the same evidence highlights HPA axis dysregulation through
\gene{NR3C1}--\gene{FKBP5}--\gene{CRH} interactions, it explains why
pharmacotherapy alone often fails: no selective serotonin reuptake inhibitor
addresses glucocorticoid receptor sensitivity or cortisol feedback dynamics.
This gap motivates the integration of psychotherapeutic approaches — notably
Dialectical Behavior Therapy, which teaches skills that functionally compensate
for the regulatory deficits these genetic liabilities produce — and
neuromodulatory interventions that target the prefrontal circuits governing
top-down emotion regulation.\cit{24}
The partially reversible epigenetic signatures at \gene{NR3C1}, \gene{FKBP5},
and \gene{BDNF} further suggest that effective treatments do not merely manage
symptoms but may rewrite the molecular record of adversity itself.
Understanding which molecular system predominates in a given patient's
dysregulation profile could, in principle, guide treatment selection toward the
modality most likely to engage that system — a vision of precision psychiatry
that the therapeutic landscape is only beginning to realize.


%% ====================================================================
%%  SECTION: THERAPEUTIC LANDSCAPE
%% ====================================================================

\sectitle{Therapeutic Landscape for Emotional Dysregulation}

The clinical management of emotional dysregulation has evolved from a narrow
reliance on pharmacotherapy into a multimodal enterprise that integrates
psychotherapy, neuromodulation, and neuropeptide-based intervention.\cit{25}
This evolution reflects a growing recognition — informed by the genetic
architecture reviewed above — that dysregulated affect arises from
interacting neural systems that no single treatment modality can
comprehensively address.
We organize the current therapeutic landscape around three axes:
skills-based psychotherapy, with Dialectical Behavior Therapy (DBT) as its
most rigorously evaluated exemplar; pharmacological approaches that target
monoamine and neuroendocrine pathways; and emerging interventions —
neuromodulation and intranasal oxytocin — that attempt to engage regulatory
circuits directly.

\subsection*{Dialectical Behavior Therapy: Evidence Across Populations}

Marsha Linehan developed DBT in the 1980s as a treatment for chronically
suicidal individuals with borderline personality disorder (BPD), grounding it in
a biosocial model that attributes emotional dysregulation to the transaction
between biological vulnerability and invalidating environments.\cit{26}
Four decades later, DBT has generated 84 registered trials on
ClinicalTrials.gov and has expanded far beyond its original indication —
a proliferation that demands careful evaluation of where the evidence is
strong, where it is preliminary, and where critical gaps remain.

The mechanistic foundation of DBT rests on four skill modules — mindfulness,
distress tolerance, emotion regulation, and interpersonal effectiveness — that
collectively target the regulatory deficits the genetic architecture section
identified.\cit{27}
Lynch and colleagues (2006) proposed that DBT achieves its effects through
three primary mechanisms of change: enhanced behavioral capabilities (patients
acquire skills they previously lacked), improved motivation (contingency
management reduces therapy-interfering behaviors), and generalization to
natural environments (telephone coaching and skills groups bridge the gap
between session and daily life).\cit{28}
Daros and Williams (2019), in a systematic review published in the
\emph{Harvard Review of Psychiatry}, examined emotion regulation strategies
in BPD more broadly and concluded that DBT's skills training component — rather
than individual therapy alone — drives the largest improvements in regulatory
capacity.\cit{29}
This finding aligns with the molecular evidence: if emotional dysregulation
reflects a deficit in prefrontal top-down control over limbic reactivity,
then skills training functions as a behavioral prosthesis that strengthens
precisely those regulatory circuits.

Recent trials have extended DBT into populations that challenge the original
biosocial model's scope.
Goldstein and colleagues (2024) published in \emph{JAMA Psychiatry} the first
randomized controlled trial of DBT for adolescents with bipolar disorder —
a population in whom emotional dysregulation is pervasive but whose mood
episodes involve neurobiological mechanisms (e.g., circadian disruption,
mitochondrial dysfunction) that differ from those in BPD.\cit{30}
The fact that DBT demonstrated efficacy in this population suggests that its
skills target transdiagnostic regulatory processes rather than
disorder-specific pathology — a conclusion consistent with the shared genetic
architecture ($r_g \approx 0.80$) between neuroticism and depression documented
above.

Bemmouna and colleagues (2025) reported in \emph{Psychotherapy and
Psychosomatics} the results of a randomized controlled trial of DBT for
autistic adults, a population characterized by alexithymia, sensory
overwhelm, and difficulties identifying and labeling emotional
states.\cit{31}
This trial is notable because autism involves differences in interoceptive
processing and social cognition that standard DBT does not explicitly address,
raising the question of whether adapted protocols that incorporate sensory
regulation and explicit emotion identification training could further improve
outcomes.

Perhaps the most innovative extension, Norman-Nott and colleagues (2025)
demonstrated in \emph{JAMA Network Open} that online DBT reduces emotional
dysregulation in individuals with chronic pain — a finding that links affective
dysregulation to nociceptive processing and suggests that the
\gene{NR3C1}--\gene{CRH} neuroendocrine axis, which mediates both stress
reactivity and pain sensitization, may represent the biological substrate
through which DBT operates in this population.\cit{32}
The online delivery format further expands DBT's scalability, addressing
longstanding concerns about treatment access in underserved settings.

\subsection*{Pharmacotherapy: Serotonergic Modulation and Its Limits}

The pharmacological treatment of emotional dysregulation has centered on
selective serotonin reuptake inhibitors (SSRIs) since their introduction in
the late 1980s.\cit{33}
Fluoxetine (PubChem CID 3386; molecular formula C\textsubscript{17}H\textsubscript{18}F\textsubscript{3}NO;
molecular weight 309.33\,Da) exemplifies the class: it inhibits \gene{SLC6A4},
thereby blocking serotonin reuptake and increasing synaptic serotonin
availability.\cit{34}
This mechanism directly engages the Reactome ``serotonin clearance'' pathway
(R-HSA-380615) and the KEGG serotonin receptor interaction pathway (map07211)
that our enrichment analysis identified as convergent nodes of genetic risk.

Yet the clinical reality of SSRI treatment exposes the limits of targeting a
single neurotransmitter system.
Remission rates for major depression hover around 30\% with first-line SSRI
monotherapy, and emotional blunting — a paradoxical reduction in both negative
\emph{and} positive affect — affects a substantial minority of
patients.\cit{35}
The genetic architecture reviewed above explains this therapeutic ceiling:
emotional dysregulation involves not only serotonergic signaling but also
HPA axis dynamics (\gene{NR3C1}, \gene{FKBP5}, \gene{CRH}), oxytocinergic
modulation (\gene{OXTR}), and neurotrophic support (\gene{BDNF}).
An SSRI addresses the first system while leaving the others untouched.

This limitation has motivated interest in pharmacological strategies that
target the neuroendocrine axis.
KEGG pathway analysis identifies oxytocin signaling (map04921) and cortisol
synthesis and secretion (map04927) as biologically plausible intervention
points.\cit{36}
\gene{FKBP5} antagonists, which would restore glucocorticoid receptor
sensitivity and normalize HPA axis negative feedback, represent a particularly
compelling target given the gene's role in the HSP90 chaperone cycle —
though no such compound has yet advanced beyond early-phase trials.\cit{37}
The classification of emotional dysregulation within the Human Phenotype
Ontology — where ``Abnormal emotional state'' (HP:0100851) subsumes 12
descendant phenotypes and ``Dysregulated negative emotional state''
(HP:0031467) subsumes 7 — underscores that what clinicians treat as a single
target actually comprises a heterogeneous family of phenotypes, each
potentially requiring a different pharmacological approach.\cit{38}

\subsection*{Neuromodulation: Engaging Prefrontal Circuitry Directly}

Where pharmacotherapy modulates neurotransmitter dynamics from below,
neuromodulation attempts to strengthen top-down regulatory circuits from above.
Neacsiu and colleagues (2021) conducted the most informative trial to date,
combining repetitive transcranial magnetic stimulation (rTMS) at 10\,Hz over
the dorsolateral prefrontal cortex (dlPFC) with cognitive restructuring — a
design that tests whether enhancing prefrontal excitability during active
skill use produces synergistic effects.\cit{39}
The results were encouraging but modest: effect sizes for emotion regulation
improvement ranged from $d = 0.20$ to $d = 0.67$, while distress reduction
showed more consistent but smaller effects ($d = 0.32$--$0.35$).

These effect sizes deserve careful interpretation.
The lower bound ($d = 0.20$) falls below conventional thresholds for clinical
significance, but the upper bound ($d = 0.67$) approaches the effect sizes
reported for DBT skills training.
The variability suggests that rTMS combined with cognitive restructuring may
benefit some individuals substantially while producing minimal effects in
others — a pattern consistent with the heterogeneous genetic architecture
described above.
If prefrontal hypoactivation represents the primary regulatory deficit for
some patients (perhaps those whose risk concentrates in monoaminergic pathways
governing cortical function) while HPA axis dysregulation predominates in
others (those with \gene{FKBP5}/\gene{NR3C1} risk variants), then
neuromodulation would be expected to show precisely this pattern of variable
response.\cit{40}

\subsection*{Intranasal Oxytocin: Neuropeptide Intervention}

The identification of \gene{OXTR} within the genetic risk architecture, and its
enrichment in both ``maternal behavior'' and ``neuroactive ligand--receptor
interaction'' pathways, provides the rationale for intranasal oxytocin as a
therapeutic strategy.\cit{41}
Moshfeghinia and colleagues (2025; PMID 41074229) conducted a comprehensive
review of fMRI studies examining intranasal oxytocin's effects on neural
activity during emotional processing.\cit{42}
Their synthesis revealed that oxytocin consistently modulates amygdala
reactivity to social and emotional stimuli, though the direction of modulation
depends on context, baseline anxiety, and — critically — sex, with stronger
anxiolytic effects in men than in women in most paradigms.

This sex-differential response highlights a broader challenge in translating
genetic findings into therapeutics.
GWAS identify risk loci across populations, but gene expression, hormonal
milieu, and receptor density vary by sex, developmental stage, and tissue
type.
The KEGG oxytocin signaling pathway (map04921) maps the molecular cascade from
receptor binding through phospholipase C activation to downstream
transcriptional effects, but this canonical pathway operates differently in
the male versus female brain, in the postpartum versus non-postpartum state,
and in the context of high versus low endogenous cortisol.\cit{43}
Effective clinical deployment of intranasal oxytocin will therefore require
stratification strategies that account for these sources of variability —
strategies that the genetic and epigenetic architecture reviewed above can, in
principle, inform.

\subsection*{Toward Integrated and Precision Approaches}

The therapeutic landscape for emotional dysregulation currently offers three
modalities that each engage a different level of the regulatory hierarchy:
DBT trains behavioral skills that strengthen prefrontal top-down control;
SSRIs modulate serotonergic neurotransmission at the synaptic level; and
neuromodulation enhances prefrontal excitability through direct cortical
stimulation.
Emerging approaches — \gene{FKBP5} antagonists, intranasal oxytocin, and
epigenetically informed interventions targeting methylation at \gene{NR3C1} and
\gene{BDNF} — promise to address the neuroendocrine and neurotrophic systems
that current treatments leave unengaged.\cit{44}

The partially reversible epigenetic signatures documented by Udalov et al.
(2026) suggest a future in which treatment response is monitored not only
through symptom questionnaires but through methylation assays at specific CpG
sites — converting molecular markers of adversity into biomarkers of
recovery.\cit{45}
miRNA panels (miR-16, miR-124, miR-132, miR-135a, miR-34c) that track limbic
remodeling could similarly serve as treatment-responsive biomarkers, enabling
clinicians to detect biological change before patients report subjective
improvement.\cit{46}

Realizing this vision requires two advances that the field has not yet achieved.
First, researchers must conduct GWAS-informed clinical trials that stratify
participants by polygenic risk scores for neuroticism, depression, and anxiety,
testing whether genetic profile predicts differential response to DBT versus
pharmacotherapy versus neuromodulation.
Second, the field must develop integrated treatment protocols that combine
modalities targeting different levels of the regulatory hierarchy —
for example, pairing rTMS-augmented DBT with SSRI pharmacotherapy and
monitoring epigenetic biomarkers to guide dose adjustment and treatment
duration.
The genetic architecture reviewed above provides the biological map;
the therapeutic landscape provides the tools.
The challenge ahead is to match the right tools to the right biological
profiles — transforming emotional dysregulation treatment from a trial-and-error
enterprise into a precision discipline.\cit{47}
